\documentclass[11pt]{article}

\usepackage[margin=1in]{geometry}
\usepackage{amsmath,amssymb}
\usepackage{bm}
\usepackage{array}
\usepackage{booktabs}
\usepackage[hidelinks]{hyperref}

\title{Part 1: RC Representations and Phase-D Architectures}
\author{}
\date{}

\begin{document}

\maketitle

We will do this as five mathematical blocks, one at a time, and \textbf{nothing in Parts 2--5 will be allowed to redefine the RC architectures established here}.

\section*{Part 1 --- RC Representations and Phase-D Architectures}

We have two RC structures:

\begin{align}
\boxed{\text{1R1C}}
\qquad\text{and}\qquad
\boxed{\text{2R2C}}
\end{align}

and four effective Phase-D modeling architectures:

\begin{align}
\boxed{\text{All-to-One}}
\end{align}

\begin{align}
\boxed{\text{Identity / IND}}
\end{align}

\begin{align}
\boxed{\text{Identity / DEP1}}
\end{align}

\begin{align}
\boxed{\text{Identity / DEP2}}.
\end{align}

For all-to-one, the Phase-D \texttt{ind/dep1/dep2} products were numerically identical, so there is only one unique all-to-one model to train. For identity, all three architectures are genuinely different.

\section*{1. Common notation}

The RestaurantFastFood identity representation has two thermal zones:

\begin{align}
i\in\{D,K\}
\end{align}

where

\begin{align}
D&=\text{Dining},\\
K&=\text{Kitchen}.
\end{align}

The all-to-one representation has one aggregate zone:

\begin{align}
A=\text{RestaurantFastFood\_All}.
\end{align}

The sampling interval is

\begin{align}
\boxed{\Delta t=300\ {\rm s}}
\end{align}

but the RC equations are formulated in continuous time.

The Phase-D \texttt{l1\_h1} data provide the measured state and forcing quantities at \(t_k\), with the temperature target at \(t_{k+1}\).

\section*{2. Heat-input definitions}

This distinction is now locked.

For any zone \(i\), the \textbf{convective} heat input is

\begin{align}
\boxed{
Q_{c,i}
=
Q_{\mathrm{ZIC},i}
+
Q_{\mathrm{Sol1},i}
+
Q_{\mathrm{AC},i}.
}
\end{align}

The \textbf{radiative} heat input is

\begin{align}
\boxed{
Q_{r,i}
=
Q_{\mathrm{ZIR},i}
+
Q_{\mathrm{Sol2},i}.
}
\end{align}

Thus

\begin{align}
Q_{c,i}
&=
Q_{zic,i}+Q_{sol1,i}+Q_{ac,i},
\\
Q_{r,i}
&=
Q_{zir,i}+Q_{sol2,i}.
\end{align}

There are no independent gain parameters multiplying these five Phase-D signals in the primary paper models.

In particular,

\begin{align}
\boxed{\alpha_{AC}=1}
\end{align}

because \texttt{QAC} has already been constructed/calibrated through Phase C.

For a structurally absent Phase-D signal, such as Kitchen solar gains in the audited identity data,

\begin{align}
Q_{\mathrm{Sol1},K}=Q_{\mathrm{Sol2},K}=0
\end{align}

by structural definition; we are not imputing an unknown solar signal.

\section*{3. 1R1C elemental zone model}

For one zone \(i\), define

\begin{align}
T_i &= \text{zone-air/effective lumped temperature},\\
T_o &= \text{outdoor temperature},\\
C_i &>0=\text{effective thermal capacitance},\\
R_{io} &>0=\text{effective zone-to-outdoor resistance}.
\end{align}

Because there is only one capacitance, the radiative input has no explicit surface/mass state into which it can be deposited.

Therefore we use the previously agreed effective radiative participation factor

\begin{align}
\boxed{
0\le \eta_{r,i}^{(1C)}\le1.
}
\end{align}

The isolated-zone equation is

\begin{align}
\boxed{
C_i\dot T_i
=
\frac{T_o-T_i}{R_{io}}
+
Q_{c,i}
+
\eta_{r,i}^{(1C)}Q_{r,i}.
}
\tag{1}
\end{align}

Expanding the heat channels,

\begin{align}
\boxed{
C_i\dot T_i
=
\frac{T_o-T_i}{R_{io}}
+
Q_{\mathrm{ZIC},i}
+
Q_{\mathrm{Sol1},i}
+
Q_{\mathrm{AC},i}
+
\eta_{r,i}^{(1C)}
\left(
Q_{\mathrm{ZIR},i}
+
Q_{\mathrm{Sol2},i}
\right).
}
\tag{2}
\end{align}

For 1R1C, \(\eta_r\) means:

\begin{quote}
effective fraction of the radiative heat represented as contributing to the single lumped thermal state.
\end{quote}

It is \textbf{not} an air/mass split because no mass state exists.

The options are

\begin{align}
\eta_r^{(1C)}=
\begin{cases}
1, & \texttt{full},\\[2mm]
0, & \texttt{zero},\\[2mm]
\eta_{\rm fixed}, & \texttt{fixed},\\[2mm]
\sigma(\rho_\eta), & \texttt{learnable},
\end{cases}
\tag{3}
\end{align}

where

\begin{align}
\sigma(z)=\frac{1}{1+e^{-z}}.
\end{align}

The main/default physical case is

\begin{align}
\boxed{\eta_r^{(1C)}=1.}
\end{align}

\section*{4. 2R2C elemental zone model}

For 2R2C, every zone has two physical states:

\begin{align}
T_{a,i}&=\text{zone-air temperature},\\
T_{m,i}&=\text{effective thermal-mass/surface temperature}.
\end{align}

The state is

\begin{align}
\boxed{
x_i=
\begin{bmatrix}
T_{a,i}\\
T_{m,i}
\end{bmatrix}.
}
\end{align}

Only \(T_{a,i}\) is measured:

\begin{align}
\boxed{
y_i=
H_ix_i,
\qquad
H_i=
\begin{bmatrix}
1&0
\end{bmatrix}.
}
\tag{4}
\end{align}

The parameters are

\begin{align}
C_{a,i}&>0,
\qquad
C_{m,i}>0,
\\
R_{ao,i}&>0,
\qquad
R_{am,i}>0.
\end{align}

Here \(\eta_r\) has a different, stronger physical interpretation:

\begin{align}
\boxed{
0\le\eta_{r,i}^{(2C)}\le1
}
\end{align}

is the fraction of radiative heat initially deposited in the effective mass state.

Therefore

\begin{align}
Q_{r\rightarrow m,i}
=
\eta_{r,i}^{(2C)}Q_{r,i},
\tag{5}
\end{align}

and

\begin{align}
Q_{r\rightarrow a,i}
=
\left(1-\eta_{r,i}^{(2C)}\right)Q_{r,i}.
\tag{6}
\end{align}

Conservation gives

\begin{align}
Q_{r\rightarrow a,i}
+
Q_{r\rightarrow m,i}
=
Q_{r,i}.
\tag{7}
\end{align}

The air equation is

\begin{align}
\boxed{
C_{a,i}\dot T_{a,i}
=
\frac{T_o-T_{a,i}}{R_{ao,i}}
+
\frac{T_{m,i}-T_{a,i}}{R_{am,i}}
+
Q_{c,i}
+
(1-\eta_{r,i}^{(2C)})Q_{r,i}.
}
\tag{8}
\end{align}

The mass equation is

\begin{align}
\boxed{
C_{m,i}\dot T_{m,i}
=
\frac{T_{a,i}-T_{m,i}}{R_{am,i}}
+
\eta_{r,i}^{(2C)}Q_{r,i}.
}
\tag{9}
\end{align}

Adding (8) and (9),

\begin{align}
\boxed{
C_{a,i}\dot T_{a,i}
+
C_{m,i}\dot T_{m,i}
=
\frac{T_o-T_{a,i}}{R_{ao,i}}
+
Q_{c,i}
+
Q_{r,i}.
}
\tag{10}
\end{align}

Notice something fundamental:

\begin{align}
R_{am,i}
\end{align}

and

\begin{align}
\eta_{r,i}
\end{align}

cancel from the total-energy equation.

That fact will become central to EBP-PINODE.

The options for 2R2C are

\begin{align}
\eta_r^{(2C)}=
\begin{cases}
1, & \texttt{mass\_only},\\
0, & \texttt{air\_only},\\
\eta_{\rm fixed}, & \texttt{fixed},\\
\sigma(\rho_\eta), & \texttt{learnable}.
\end{cases}
\tag{11}
\end{align}

The main/default model is

\begin{align}
\boxed{\eta_r^{(2C)}=1.}
\end{align}

\section*{5. Architecture A --- All-to-One}

The entire Dining+Kitchen building representation is reduced to one aggregate thermal zone \(A\).

Phase D supplies

\begin{align}
T_A,
\quad
Q_{AC,A},
\quad
T_o,
\end{align}

\begin{align}
Q_{ZIC,A},
\quad
Q_{ZIR,A},
\quad
Q_{Sol1,A},
\quad
Q_{Sol2,A}.
\end{align}

Define

\begin{align}
Q_{c,A}
=
Q_{ZIC,A}+Q_{Sol1,A}+Q_{AC,A},
\tag{12}
\end{align}

\begin{align}
Q_{r,A}
=
Q_{ZIR,A}+Q_{Sol2,A}.
\tag{13}
\end{align}

There is no inter-zone resistance because only one aggregate state remains.

\subsection*{5.1 All-to-One 1R1C}

State:

\begin{align}
\boxed{x_A=T_A.}
\end{align}

Parameters:

\begin{align}
\boxed{
\theta_{A,1C}
=
\{R_{Ao},C_A,\eta_{r,A}\}.
}
\end{align}

Dynamics:

\begin{align}
\boxed{
C_A\dot T_A
=
\frac{T_o-T_A}{R_{Ao}}
+
Q_{c,A}
+
\eta_{r,A}Q_{r,A}.
}
\tag{14}
\end{align}

Equivalently,

\begin{align}
\dot T_A
=
\frac{1}{C_A}
\left[
\frac{T_o-T_A}{R_{Ao}}
+
Q_{c,A}
+
\eta_{r,A}Q_{r,A}
\right].
\tag{15}
\end{align}

The measured output is

\begin{align}
y_A=T_A.
\end{align}

\section*{6. All-to-One 2R2C}

State:

\begin{align}
\boxed{
x_A=
\begin{bmatrix}
T_{a,A}\\
T_{m,A}
\end{bmatrix}.
}
\end{align}

Measured output:

\begin{align}
y_A=T_{a,A}.
\end{align}

Parameters:

\begin{align}
\boxed{
\theta_{A,2C}
=
\{
R_{Ao},
R_{Am},
C_{a,A},
C_{m,A},
\eta_{r,A}
\}.
}
\end{align}

Dynamics:

\begin{align}
\boxed{
C_{a,A}\dot T_{a,A}
=
\frac{T_o-T_{a,A}}{R_{Ao}}
+
\frac{T_{m,A}-T_{a,A}}{R_{Am}}
+
Q_{c,A}
+
(1-\eta_{r,A})Q_{r,A}.
}
\tag{16}
\end{align}

\begin{align}
\boxed{
C_{m,A}\dot T_{m,A}
=
\frac{T_{a,A}-T_{m,A}}{R_{Am}}
+
\eta_{r,A}Q_{r,A}.
}
\tag{17}
\end{align}

Only

\begin{align}
T_{m,A}(t_0)
\end{align}

is hidden.

\section*{7. Architecture B --- Identity / IND}

Phase D's IND architecture treats Dining and Kitchen as \textbf{two independent systems}.

This means no inter-zone state is supplied to the other zone's model.

Therefore

\begin{align}
\boxed{R_{DK}\text{ does not exist in IND}.}
\end{align}

This is an architectural constraint, not merely setting some learned coupling to a large value.

\section*{8. Identity / IND --- 1R1C}

Dining:

\begin{align}
\boxed{
C_D\dot T_D
=
\frac{T_o-T_D}{R_{Do}}
+
Q_{c,D}
+
\eta_{r,D}Q_{r,D}.
}
\tag{18}
\end{align}

Kitchen:

\begin{align}
\boxed{
C_K\dot T_K
=
\frac{T_o-T_K}{R_{Ko}}
+
Q_{c,K}
+
\eta_{r,K}Q_{r,K}.
}
\tag{19}
\end{align}

These are two independent models:

\begin{align}
f_D(T_D,u_D,d_D)
\end{align}

and

\begin{align}
f_K(T_K,u_K,d_K).
\end{align}

Their state dimensions are each

\begin{align}
n_x=1.
\end{align}

The complete building simulator may still execute both simultaneously, but mathematically they do not communicate.

For Kitchen,

\begin{align}
Q_{Sol1,K}=Q_{Sol2,K}=0,
\end{align}

so

\begin{align}
Q_{c,K}=Q_{ZIC,K}+Q_{AC,K},
\tag{20}
\end{align}

\begin{align}
Q_{r,K}=Q_{ZIR,K}.
\tag{21}
\end{align}

We retain the generic equations so the architecture does not need special code other than the absence mask.

\section*{9. Identity / IND --- 2R2C}

Dining state:

\begin{align}
x_D=
\begin{bmatrix}
T_{a,D}\\T_{m,D}
\end{bmatrix}.
\end{align}

Kitchen state:

\begin{align}
x_K=
\begin{bmatrix}
T_{a,K}\\T_{m,K}
\end{bmatrix}.
\end{align}

Dining:

\begin{align}
\boxed{
C_{a,D}\dot T_{a,D}
=
\frac{T_o-T_{a,D}}{R_{Do}}
+
\frac{T_{m,D}-T_{a,D}}{R_{Dm}}
+
Q_{c,D}
+
(1-\eta_{r,D})Q_{r,D}.
}
\tag{22}
\end{align}

\begin{align}
\boxed{
C_{m,D}\dot T_{m,D}
=
\frac{T_{a,D}-T_{m,D}}{R_{Dm}}
+
\eta_{r,D}Q_{r,D}.
}
\tag{23}
\end{align}

Kitchen:

\begin{align}
\boxed{
C_{a,K}\dot T_{a,K}
=
\frac{T_o-T_{a,K}}{R_{Ko}}
+
\frac{T_{m,K}-T_{a,K}}{R_{Km}}
+
Q_{c,K}
+
(1-\eta_{r,K})Q_{r,K}.
}
\tag{24}
\end{align}

\begin{align}
\boxed{
C_{m,K}\dot T_{m,K}
=
\frac{T_{a,K}-T_{m,K}}{R_{Km}}
+
\eta_{r,K}Q_{r,K}.
}
\tag{25}
\end{align}

Again,

\begin{align}
R_{DK}
\end{align}

is absent.

\section*{10. Architecture C --- Identity / DEP1}

This is our primary physically coupled identity architecture.

Phase D supplies both zone temperatures:

\begin{align}
T_D,\quad T_K,
\end{align}

both zone QAC values:

\begin{align}
Q_{AC,D},\quad Q_{AC,K},
\end{align}

and zone-specific disturbances.

We therefore introduce one physical inter-zone resistance:

\begin{align}
\boxed{R_{DK}>0.}
\end{align}

We use reciprocal coupling:

\begin{align}
Q_{D\rightarrow K}
=
\frac{T_D-T_K}{R_{DK}},
\tag{26}
\end{align}

and therefore

\begin{align}
Q_{K\rightarrow D}
=
-\,
Q_{D\rightarrow K}
=
\frac{T_K-T_D}{R_{DK}}.
\tag{27}
\end{align}

Thus inter-zone exchange conserves energy:

\begin{align}
Q_{D\rightarrow K}+Q_{K\rightarrow D}=0.
\tag{28}
\end{align}

This reciprocal/shared \(R_{DK}\) assumption is important for both physical interpretation and the later EBP projection.

\section*{11. Identity / DEP1 --- coupled 1R1C}

State:

\begin{align}
\boxed{
x=
\begin{bmatrix}
T_D\\
T_K
\end{bmatrix}.
}
\end{align}

Mass/capacitance matrix:

\begin{align}
\boxed{
M_{1C}
=
\begin{bmatrix}
C_D&0\\
0&C_K
\end{bmatrix}.
}
\tag{29}
\end{align}

Dynamics:

\begin{align}
\boxed{
C_D\dot T_D
=
\frac{T_o-T_D}{R_{Do}}
+
\frac{T_K-T_D}{R_{DK}}
+
Q_{c,D}
+
\eta_{r,D}Q_{r,D}.
}
\tag{30}
\end{align}

\begin{align}
\boxed{
C_K\dot T_K
=
\frac{T_o-T_K}{R_{Ko}}
+
\frac{T_D-T_K}{R_{DK}}
+
Q_{c,K}
+
\eta_{r,K}Q_{r,K}.
}
\tag{31}
\end{align}

Define

\begin{align}
b_D=
\frac{T_o-T_D}{R_{Do}}
+
\frac{T_K-T_D}{R_{DK}}
+
Q_{c,D}
+
\eta_{r,D}Q_{r,D},
\tag{32}
\end{align}

\begin{align}
b_K=
\frac{T_o-T_K}{R_{Ko}}
+
\frac{T_D-T_K}{R_{DK}}
+
Q_{c,K}
+
\eta_{r,K}Q_{r,K}.
\tag{33}
\end{align}

Then

\begin{align}
\boxed{
M_{1C}\dot x=
\begin{bmatrix}
b_D\\b_K
\end{bmatrix}.
}
\tag{34}
\end{align}

or

\begin{align}
\boxed{
\dot x
=
M_{1C}^{-1}
b(x,u,d;\theta).
}
\tag{35}
\end{align}

\section*{12. Identity / DEP1 --- coupled 2R2C}

State ordering:

\begin{align}
\boxed{
x=
\begin{bmatrix}
T_{a,D}\\
T_{m,D}\\
T_{a,K}\\
T_{m,K}
\end{bmatrix}.
}
\tag{36}
\end{align}

Measured output:

\begin{align}
\boxed{
y=
\begin{bmatrix}
T_{a,D}\\
T_{a,K}
\end{bmatrix}
=
Hx,
}
\end{align}

with

\begin{align}
\boxed{
H=
\begin{bmatrix}
1&0&0&0\\
0&0&1&0
\end{bmatrix}.
}
\tag{37}
\end{align}

The two hidden states are

\begin{align}
T_{m,D},\qquad T_{m,K}.
\end{align}

Capacitance matrix:

\begin{align}
\boxed{
M_{2C}
=
\operatorname{diag}
(C_{a,D},C_{m,D},C_{a,K},C_{m,K}).
}
\tag{38}
\end{align}

Dining air:

\begin{align}
\boxed{
\begin{aligned}
C_{a,D}\dot T_{a,D}
={}&
\frac{T_o-T_{a,D}}{R_{Do}}
+
\frac{T_{m,D}-T_{a,D}}{R_{Dm}}\\
&+
\frac{T_{a,K}-T_{a,D}}{R_{DK}}
+
Q_{c,D}
+
(1-\eta_{r,D})Q_{r,D}.
\end{aligned}
}
\tag{39}
\end{align}

Dining mass:

\begin{align}
\boxed{
C_{m,D}\dot T_{m,D}
=
\frac{T_{a,D}-T_{m,D}}{R_{Dm}}
+
\eta_{r,D}Q_{r,D}.
}
\tag{40}
\end{align}

Kitchen air:

\begin{align}
\boxed{
\begin{aligned}
C_{a,K}\dot T_{a,K}
={}&
\frac{T_o-T_{a,K}}{R_{Ko}}
+
\frac{T_{m,K}-T_{a,K}}{R_{Km}}\\
&+
\frac{T_{a,D}-T_{a,K}}{R_{DK}}
+
Q_{c,K}
+
(1-\eta_{r,K})Q_{r,K}.
\end{aligned}
}
\tag{41}
\end{align}

Kitchen mass:

\begin{align}
\boxed{
C_{m,K}\dot T_{m,K}
=
\frac{T_{a,K}-T_{m,K}}{R_{Km}}
+
\eta_{r,K}Q_{r,K}.
}
\tag{42}
\end{align}

This is the most spatially detailed physical RC representation in the paper.

\section*{13. Total-energy equations for DEP1 2R2C}

Add Dining equations (39)+(40):

\begin{align}
\boxed{
\begin{aligned}
C_{a,D}\dot T_{a,D}
+
C_{m,D}\dot T_{m,D}
={}&
\frac{T_o-T_{a,D}}{R_{Do}}\\
&+
\frac{T_{a,K}-T_{a,D}}{R_{DK}}
+
Q_{c,D}+Q_{r,D}.
\end{aligned}
}
\tag{43}
\end{align}

Likewise,

\begin{align}
\boxed{
\begin{aligned}
C_{a,K}\dot T_{a,K}
+
C_{m,K}\dot T_{m,K}
={}&
\frac{T_o-T_{a,K}}{R_{Ko}}\\
&+
\frac{T_{a,D}-T_{a,K}}{R_{DK}}
+
Q_{c,K}+Q_{r,K}.
\end{aligned}
}
\tag{44}
\end{align}

Again,

\begin{align}
\eta_r
\end{align}

and the air-mass conductances disappear from these total balances.

If we add the two zone totals,

\begin{align}
&C_{a,D}\dot T_{a,D}
+C_{m,D}\dot T_{m,D}
+
C_{a,K}\dot T_{a,K}
+C_{m,K}\dot T_{m,K},
\end{align}

then the inter-zone exchange also cancels:

\begin{align}
\frac{T_K-T_D}{R_{DK}}
+
\frac{T_D-T_K}{R_{DK}}
=0.
\end{align}

Therefore the complete two-zone building balance is

\begin{align}
\boxed{
\begin{aligned}
\dot E_{\rm building}
={}&
\frac{T_o-T_{a,D}}{R_{Do}}
+
\frac{T_o-T_{a,K}}{R_{Ko}}\\
&+
Q_{c,D}+Q_{r,D}
+
Q_{c,K}+Q_{r,K}.
\end{aligned}
}
\tag{45}
\end{align}

This will be valuable later for both our aggregation theorem and EBP derivation.

\section*{14. Architecture D --- Identity / DEP2}

DEP2 requires the most care.

Phase D gives us:

\subsection*{Zone-specific states}

\begin{align}
T_D,\quad T_K.
\end{align}

\subsection*{Zone-specific HVAC thermal controls}

\begin{align}
Q_{AC,D},
\quad
Q_{AC,K}.
\end{align}

\subsection*{Common all-to-one non-HVAC disturbances}

\begin{align}
Q_{ZIC,A},
\quad
Q_{Sol1,A},
\end{align}

\begin{align}
Q_{ZIR,A},
\quad
Q_{Sol2,A}.
\end{align}

So define the aggregate non-HVAC convective signal

\begin{align}
\boxed{
\bar Q_c^{\,nh}
=
Q_{ZIC,A}+Q_{Sol1,A}.
}
\tag{46}
\end{align}

and aggregate radiative signal

\begin{align}
\boxed{
\bar Q_r
=
Q_{ZIR,A}+Q_{Sol2,A}.
}
\tag{47}
\end{align}

I use the overbar deliberately because these came from the equal-weight all-to-one aggregation.

The numerical audit is consistent with those signals behaving as equal-weight aggregate quantities rather than zone-resolved individual heat gains.

Therefore we must not write

\begin{align}
Q_D=\bar Q
\end{align}

and

\begin{align}
Q_K=\bar Q
\end{align}

without qualification.

That would duplicate the aggregate heat.

\section*{15. DEP2 effective disturbance allocation}

We introduce separate allocation factors for the convective non-HVAC and radiative disturbances:

\begin{align}
\lambda_{c,D},
\quad
\lambda_{c,K},
\end{align}

\begin{align}
\lambda_{r,D},
\quad
\lambda_{r,K}.
\end{align}

Since the all-to-one signals are equal-weight two-zone aggregate quantities, impose

\begin{align}
\boxed{
\lambda_{c,D}+\lambda_{c,K}=2,
}
\tag{48}
\end{align}

\begin{align}
\boxed{
\lambda_{r,D}+\lambda_{r,K}=2.
}
\tag{49}
\end{align}

Then

\begin{align}
\boxed{
Q_{c,D}^{DEP2}
=
Q_{AC,D}
+
\lambda_{c,D}\bar Q_c^{nh},
}
\tag{50}
\end{align}

\begin{align}
\boxed{
Q_{c,K}^{DEP2}
=
Q_{AC,K}
+
\lambda_{c,K}\bar Q_c^{nh}.
}
\tag{51}
\end{align}

Radiative:

\begin{align}
\boxed{
Q_{r,D}^{DEP2}
=
\lambda_{r,D}\bar Q_r,
}
\tag{52}
\end{align}

\begin{align}
\boxed{
Q_{r,K}^{DEP2}
=
\lambda_{r,K}\bar Q_r.
}
\tag{53}
\end{align}

This preserves

\begin{align}
\frac{
\lambda_{c,D}\bar Q_c^{nh}
+
\lambda_{c,K}\bar Q_c^{nh}
}{2}
=
\bar Q_c^{nh},
\end{align}

and likewise for radiation.

\section*{16. Convenient constrained parameterization}

With two zones, we only need one unconstrained parameter for each allocation family.

Let

\begin{align}
\alpha_c\in\mathbb R,
\qquad
\alpha_r\in\mathbb R.
\end{align}

Set

\begin{align}
\boxed{
\lambda_{c,D}=2\sigma(\alpha_c)
}
\tag{54}
\end{align}

and

\begin{align}
\boxed{
\lambda_{c,K}=2[1-\sigma(\alpha_c)].
}
\tag{55}
\end{align}

Then automatically

\begin{align}
\lambda_{c,D}>0,
\qquad
\lambda_{c,K}>0,
\end{align}

and

\begin{align}
\lambda_{c,D}+\lambda_{c,K}=2.
\end{align}

Similarly,

\begin{align}
\boxed{
\lambda_{r,D}=2\sigma(\alpha_r),
}
\tag{56}
\end{align}

\begin{align}
\boxed{
\lambda_{r,K}=2[1-\sigma(\alpha_r)].
}
\tag{57}
\end{align}

Equal allocation corresponds to

\begin{align}
\alpha_c=\alpha_r=0
\end{align}

because

\begin{align}
\sigma(0)=\frac12
\end{align}

and therefore

\begin{align}
\lambda_D=\lambda_K=1.
\end{align}

This is a natural initialization.

These are \textbf{DEP2 information-allocation parameters}, not generic heat multipliers.

That distinction should remain explicit in the paper.

\section*{17. Identity / DEP2 --- 1R1C}

The state structure is exactly the same as DEP1:

\begin{align}
x=
\begin{bmatrix}
T_D\\T_K
\end{bmatrix}.
\end{align}

The coupling is also the same:

\begin{align}
R_{DK}>0.
\end{align}

What changes is only the non-HVAC disturbance information supplied to each zone.

Dining:

\begin{align}
\boxed{
\begin{aligned}
C_D\dot T_D
={}&
\frac{T_o-T_D}{R_{Do}}
+
\frac{T_K-T_D}{R_{DK}}\\
&+
Q_{AC,D}
+
\lambda_{c,D}\bar Q_c^{nh}\\
&+
\eta_{r,D}
\lambda_{r,D}\bar Q_r.
\end{aligned}
}
\tag{58}
\end{align}

Kitchen:

\begin{align}
\boxed{
\begin{aligned}
C_K\dot T_K
={}&
\frac{T_o-T_K}{R_{Ko}}
+
\frac{T_D-T_K}{R_{DK}}\\
&+
Q_{AC,K}
+
\lambda_{c,K}\bar Q_c^{nh}\\
&+
\eta_{r,K}
\lambda_{r,K}\bar Q_r.
\end{aligned}
}
\tag{59}
\end{align}

\section*{18. Identity / DEP2 --- 2R2C}

Same four-state structure as DEP1:

\begin{align}
x=
\begin{bmatrix}
T_{a,D}\\T_{m,D}\\T_{a,K}\\T_{m,K}
\end{bmatrix}.
\end{align}

Dining air:

\begin{align}
\boxed{
\begin{aligned}
C_{a,D}\dot T_{a,D}
={}&
\frac{T_o-T_{a,D}}{R_{Do}}
+
\frac{T_{m,D}-T_{a,D}}{R_{Dm}}\\
&+
\frac{T_{a,K}-T_{a,D}}{R_{DK}}
+
Q_{AC,D}
+
\lambda_{c,D}\bar Q_c^{nh}\\
&+
(1-\eta_{r,D})
\lambda_{r,D}\bar Q_r.
\end{aligned}
}
\tag{60}
\end{align}

Dining mass:

\begin{align}
\boxed{
C_{m,D}\dot T_{m,D}
=
\frac{T_{a,D}-T_{m,D}}{R_{Dm}}
+
\eta_{r,D}
\lambda_{r,D}\bar Q_r.
}
\tag{61}
\end{align}

Kitchen air:

\begin{align}
\boxed{
\begin{aligned}
C_{a,K}\dot T_{a,K}
={}&
\frac{T_o-T_{a,K}}{R_{Ko}}
+
\frac{T_{m,K}-T_{a,K}}{R_{Km}}\\
&+
\frac{T_{a,D}-T_{a,K}}{R_{DK}}
+
Q_{AC,K}
+
\lambda_{c,K}\bar Q_c^{nh}\\
&+
(1-\eta_{r,K})
\lambda_{r,K}\bar Q_r.
\end{aligned}
}
\tag{62}
\end{align}

Kitchen mass:

\begin{align}
\boxed{
C_{m,K}\dot T_{m,K}
=
\frac{T_{a,K}-T_{m,K}}{R_{Km}}
+
\eta_{r,K}
\lambda_{r,K}\bar Q_r.
}
\tag{63}
\end{align}

\section*{19. DEP2 total-energy balances}

Dining:

\begin{align}
\boxed{
\begin{aligned}
C_{a,D}\dot T_{a,D}
+
C_{m,D}\dot T_{m,D}
={}&
\frac{T_o-T_{a,D}}{R_{Do}}
+
\frac{T_{a,K}-T_{a,D}}{R_{DK}}\\
&+
Q_{AC,D}
+
\lambda_{c,D}\bar Q_c^{nh}
+
\lambda_{r,D}\bar Q_r.
\end{aligned}
}
\tag{64}
\end{align}

Kitchen:

\begin{align}
\boxed{
\begin{aligned}
C_{a,K}\dot T_{a,K}
+
C_{m,K}\dot T_{m,K}
={}&
\frac{T_o-T_{a,K}}{R_{Ko}}
+
\frac{T_{a,D}-T_{a,K}}{R_{DK}}\\
&+
Q_{AC,K}
+
\lambda_{c,K}\bar Q_c^{nh}
+
\lambda_{r,K}\bar Q_r.
\end{aligned}
}
\tag{65}
\end{align}

Again,

\begin{align}
\eta_r
\end{align}

cancels from the total balance.

\section*{20. What exactly is observed versus hidden?}

This needs to remain fixed for all four learning methods.

\subsection*{1R1C}

Everything in the physical state is observed:

\begin{align}
x=y=T_z.
\end{align}

No latent thermal state exists.

\subsection*{2R2C}

Only zone air is observed.

For one zone:

\begin{align}
x=
\begin{bmatrix}
T_a\\T_m
\end{bmatrix},
\qquad
y=T_a.
\end{align}

For coupled identity:

\begin{align}
x=
\begin{bmatrix}
T_{a,D}\\
T_{m,D}\\
T_{a,K}\\
T_{m,K}
\end{bmatrix}.
\end{align}

\begin{align}
y=
\begin{bmatrix}
T_{a,D}\\
T_{a,K}
\end{bmatrix}.
\end{align}

Thus the unobserved variables are exactly

\begin{align}
T_{m,D},T_{m,K}.
\end{align}

No learning method gets to invent a latent coordinate system that replaces the physical states.

That becomes particularly important for NODE/PINODE/EBP.

\section*{21. Exact Phase-D architecture map}

This is the interface I would freeze for the later TeX.

\begin{center}
\begin{tabular}{|l|l|l|l|l|}
\hline
Architecture & State information & QAC & Non-HVAC disturbances & Inter-zone coupling\\
\hline
All-to-One & aggregate & aggregate & aggregate & none\\
Identity/IND & zone-specific, separate & zone-specific & zone-specific & none\\
Identity/DEP1 & joint zone-specific & zone-specific & zone-specific & yes\\
Identity/DEP2 & joint zone-specific & zone-specific & all-to-one aggregate & yes + allocation map\\
\hline
\end{tabular}
\end{center}

And by RC type:

\begin{center}
\begin{tabular}{|l|c|c|}
\hline
Architecture & 1R1C states & 2R2C states\\
\hline
All-to-One & \(T_A\) & \(T_{a,A},T_{m,A}\)\\
IND Dining & \(T_D\) & \(T_{a,D},T_{m,D}\)\\
IND Kitchen & \(T_K\) & \(T_{a,K},T_{m,K}\)\\
DEP1 & \(T_D,T_K\) & \(T_{a,D},T_{m,D},T_{a,K},T_{m,K}\)\\
DEP2 & \(T_D,T_K\) & \(T_{a,D},T_{m,D},T_{a,K},T_{m,K}\)\\
\hline
\end{tabular}
\end{center}

\section*{22. Physical parameter sets}

This is useful because Parts 2, 4 and 5 will estimate these.

\subsection*{All-to-One 1R1C}

\begin{align}
\boxed{
\theta=
\{R_{Ao},C_A,\eta_{r,A}\}.
}
\end{align}

\subsection*{All-to-One 2R2C}

\begin{align}
\boxed{
\theta=
\{
R_{Ao},
R_{Am},
C_{a,A},
C_{m,A},
\eta_{r,A}
\}.
}
\end{align}

\subsection*{Identity/IND 1R1C}

Dining:

\begin{align}
\theta_D=
\{R_{Do},C_D,\eta_{r,D}\}.
\end{align}

Kitchen:

\begin{align}
\theta_K=
\{R_{Ko},C_K,\eta_{r,K}\}.
\end{align}

\subsection*{Identity/IND 2R2C}

\begin{align}
\theta_D=
\{
R_{Do},R_{Dm},
C_{a,D},C_{m,D},
\eta_{r,D}
\},
\end{align}

and similarly for Kitchen.

\subsection*{Identity/DEP1 1R1C}

\begin{align}
\boxed{
\theta=
\{
R_{Do},R_{Ko},R_{DK},
C_D,C_K,
\eta_{r,D},\eta_{r,K}
\}.
}
\end{align}

\subsection*{Identity/DEP1 2R2C}

\begin{align}
\boxed{
\begin{aligned}
\theta=\{&
R_{Do},R_{Ko},R_{DK},
R_{Dm},R_{Km},\\
&
C_{a,D},C_{m,D},
C_{a,K},C_{m,K},\\
&
\eta_{r,D},\eta_{r,K}
\}.
\end{aligned}
}
\end{align}

\subsection*{DEP2}

Add

\begin{align}
\boxed{\alpha_c,\alpha_r}
\end{align}

or equivalently the constrained

\begin{align}
\lambda_{c,D},\lambda_{c,K},
\lambda_{r,D},\lambda_{r,K}.
\end{align}

\section*{23. Positivity constraints}

Every resistance and capacitance must satisfy

\begin{align}
R>0,\qquad C>0.
\end{align}

For trainable models I recommend parameterizing them as

\begin{align}
\boxed{
R=R_{\min}+\operatorname{softplus}(\rho_R)
}
\tag{66}
\end{align}

and

\begin{align}
\boxed{
C=C_{\min}+\operatorname{softplus}(\rho_C).
}
\tag{67}
\end{align}

This is better than optimizing \(R,C\) directly and clipping afterward because positivity is maintained throughout differentiation.

Likewise,

\begin{align}
\eta_r=\sigma(\rho_\eta)
\end{align}

when learnable.

\section*{24. Important modeling decision: same RC architecture across all four methods}

This is critical for experimental fairness.

Inverse PINN-RC, NODE, base PINODE and EBP-PINODE will \textbf{not} be allowed to change what:

\begin{itemize}
\item all-to-one means,
\item IND means,
\item DEP1 means,
\item DEP2 means,
\item 1R1C means,
\item 2R2C means,
\item convective heat means,
\item radiative heat means,
\item \(\eta_r\) means,
\item inter-zone coupling means.
\end{itemize}

The methods may differ in \textbf{how the vector field is learned}, but the data/information architecture remains identical.

That gives us a clean scientific comparison.

\section*{25. Do we have all required thermal training data from Phase D?}

For these RC representations:

\begin{align}
\boxed{\text{Yes.}}
\end{align}

The audited \texttt{ML\_SciML / l1\_h1} Phase-D products supply:

\begin{itemize}
\item observed zone temperature;
\item one-step target zone temperature;
\item outdoor temperature;
\item QAC;
\item ZIC;
\item ZIR;
\item Sol1 where applicable;
\item Sol2 where applicable;
\item train/validation/test membership;
\item timestamps;
\item IND/DEP1/DEP2 information architecture.
\end{itemize}

The only quantities \textbf{not measured in Phase D} are the 2R2C hidden mass temperatures

\begin{align}
T_m,
\end{align}

but those are intentionally latent model states, not missing training inputs. Their initialization/estimation belongs to Parts 2--5.

DEP2's zone-level non-HVAC disturbances are likewise deliberately absent by construction, which is why we have introduced the explicit constrained allocation map rather than pretending those measurements exist.

So we do \textbf{not} need another upstream dataset to train the thermal models themselves.

Phase B/C becomes necessary later for Sim3 and MPC because of the

\begin{align}
\dot m\rightarrow Q_{AC}
\end{align}

and

\begin{align}
Q_{AC}\rightarrow P_{HVAC}
\end{align}

mappings---not because Phase D is insufficient for training the thermal dynamics.

\section*{Part 1 conclusion}

I recommend we now treat equations (1)--(67) above as the candidate locked RC architecture.

The key decisions are:

\begin{align}
\boxed{
Q_c=Q_{ZIC}+Q_{Sol1}+Q_{AC}
}
\end{align}

\begin{align}
\boxed{
Q_r=Q_{ZIR}+Q_{Sol2}
}
\end{align}

with no individual gain coefficients;

\begin{align}
\boxed{
\eta_r^{1C}
=
\text{effective radiative participation}
}
\end{align}

while

\begin{align}
\boxed{
\eta_r^{2C}
=
\text{radiative mass-allocation fraction};
}
\end{align}

IND has \textbf{no inter-zone resistance};

DEP1 uses a reciprocal physical

\begin{align}
R_{DK};
\end{align}

and DEP2 uses the same coupled thermal structure but adds constrained effective allocation parameters

\begin{align}
\boxed{\lambda_c,\lambda_r}
\end{align}

for the aggregate non-HVAC disturbances.

\end{document}
